% !TeX root = ../main.tex
\section{Discussion}
\label{sec:discussion}

\subsection{Principal Accuracy Bottlenecks}

The combined ResNet, ConvNeXt, and Tiny-ViT analyses indicate that the dominant limitations under the canonical deployment regime (standard 4-bit quantization with 5\% C2C and 10\% D2D variability) are not the ones most often emphasized at the device level. Once hardware-aware training (HAT) and scale recovery are available, standard cycle-to-cycle (C2C) and device-to-device (D2D) variability do not dominate the error budget. Three constraints dominate instead.

A quantitative decomposition supports this ordering. Using the first-order Sobol definition in Eq.~\ref{eq:sobol-first-order}, the $7 \times 9$ D2D--ADC grid shows that, over the full parameter space, ADC precision accounts for 97.6\% of accuracy variance ($S_{\mathrm{ADC}}=0.98$), while D2D contributes only 1.8\%. However, conditioning on the operational envelope ($\geq$6-bit ADC, $\sigma_{\mathrm{D2D}} \leq 15\%$), the hierarchy inverts: D2D variability accounts for 92.2\% of residual variance ($S_{\mathrm{D2D}}=0.92$). This two-phase structure has a direct hardware implication: designers should first secure 6-bit readout, then invest the remaining error budget in reducing device-to-device mismatch.

The first is ADC resolution. The transition around 6-bit ADC resolution marks a clear change in model behavior: below this point transformer inference degrades substantially, indicating that improvements in conductance control remain gated by readout precision in the present mixed-signal stack.

The second is hardware-instance overfitting. The collapse of V4 under fresh-instance transfer (10.00\% accuracy) shows that standard HAT readily fits a single fixed D2D realization. Replacing the fixed-mask objective in Eq.~\ref{eq:hat-standard} with the mismatch-distributed objective in Eq.~\ref{eq:hat-ensemble} raises fresh-instance accuracy to \(86.37\pm1.54\%\), indicating that multi-instance-aware training can mitigate this transfer failure mode. Mechanistically, Ensemble HAT changes the fixed spatial mismatch map itself across epochs rather than merely perturbing one nominal instance.

A residual gap to the 98.06\% digital baseline persists under matched deployment conditions. Part of this difference may reflect regularization from per-epoch D2D resampling rather than analog noise alone. Capacity-aware Ensemble HAT schedules remain an open optimization problem.

The third is the scale-masking regime observed in V2. The high accuracy of V2 under standard uniform noise is a conditional property of the canonical simulator configuration rather than a universal form of robustness. When the recovered-weight law in Eq.~\ref{eq:scale-recovery} maps analog perturbations below the 4-bit quantization step, many realizations remain in the same recovered weight bin after calibration. This protection depends on accurate post-array gain calibration and disappears once the noise law becomes state-dependent, as shown by the proportional-noise experiments.

\subsection{Transformer Sensitivity to Non-Idealities}

Comparing the dual testbeds—ConvNeXt (CNN) and Tiny-ViT (Transformer)—reveals that the Transformer architecture is more fragile under the present physical stress tests. Because the two models serve complementary roles (ConvNeXt from scratch versus Tiny-ViT from pre-trained fine-tuning), this result should be read as evidence of a strong architecture-and-training interaction rather than as a universal ranking of all Transformers against all CNNs.

The proportional-noise experiments illustrate this most clearly. The collapse of V4 under proportional noise (10.00\% accuracy), contrasted with the relative stability of ConvNeXt under moderate nonlinearities, indicates that Transformer inference depends on the precision of token projections. Because attention maps arise from global interactions, perturbations at a small number of high-conductance devices can disrupt the spatial weighting of the entire block.

\paragraph{Why the transformer amplifies front-end distortion.}
The larger accuracy drop in V6 than in ResNet R4 reflects a structural difference between the two architectures. In a CNN, local convolution and pooling average perturbations spatially, and ReLU saturation limits the propagation of intensity-dependent scaling errors. In a ViT, patch embeddings are fed directly into global self-attention: a sublinear photoresponse alters the token embedding, and the softmax exponentiates similarity changes, redistributing attention weights across the entire sequence without spatial averaging. The detailed mechanism and quantitative support are discussed in the front-end theory note of the Supplementary Information.

Robustness is highly regime-matched. Tiny-ViT can recover strongly when trained and evaluated under the same proportional-noise law, yet the resulting checkpoint does not transfer back to the canonical uniform-noise regime. In this framework, robustness is distribution-specific rather than noise-invariant.

The framework remains reasonably resilient overall. A single-run compound stress test combining all modeled non-idealities (2\% C2C, 3\% D2D, 6-bit ADC, 1\% IR drop, 1\% sneak path, and $1{,}000$ s retention) maintained an 89.61\% accuracy for Tiny-ViT, suggesting that the mixed-precision architecture with Ensemble HAT can tolerate moderate cumulative deployment stress within the present model.

\subsection{Task Complexity and Data Starvation}

HAT recovery is most valuable on CIFAR-100, where decision boundaries are fine, whereas CIFAR-10 is simple enough that noisy standard models remain functional. Flowers-102 marks the low-data extreme (Supplementary Table~\SuppTableFlowers).

\subsection{The "Analog Ceiling" and Energy Efficiency}

A first-order energy model projects roughly an order-of-magnitude dense-projection energy reduction versus FP32 (approximately 11×), reduced to approximately 10× after 10--50\% routing overhead. Digital attention still dominates total energy (roughly 60\%). These are system-level upper bounds under placeholder constants, not chip-predictive estimates.

\subsection{Limitations}
\label{subsec:limitations}

Energy parameters are placeholders; parasitic ranges derive from ReRAM literature \citep{fastirdrop2025,iconniv2025} until measured organic-array data become available. Three circuit-level phenomena are not yet modeled explicitly: spatial IR drop (currently a 1\% scalar placeholder), sneak-path currents (likewise), and temperature-dependent shifts in photoresponse and mismatch variance. Heavy-tailed conductance distributions are also absent. Spatially correlated D2D is evaluated in Supplementary Note~SX.Z: under separable AR(1) perturbation, moderate correlation ($\rho=0.3$) degrades Ensemble HAT fresh-instance accuracy by $1.76$~pp versus the matched i.i.d. baseline ($84.57\pm2.39$\% vs.\ $86.33\pm1.61$\%), and stronger correlation ($\rho=0.5$) degrades it by $4.20$~pp ($82.12\pm3.95$\%). Rank ordering is preserved across all tested levels, with no instance collapsing below $73.7$\% (\texttt{fresh\_instance\_eval\_v4\_ensemble\_correlated\_d2d.json}). Scale-masking depends on ideal calibrated rescaling; a hook-based check shows that moderate ADC calibration errors change V4 accuracy by only 0.18~pp. The energy model further assumes ideal single-shot programming and therefore excludes iterative write-verify overhead for 4-bit conductance states. The $NL=2.0$ limit reflects the present gradient-scaling surrogate, not a material bound. Group-wise ablation confirms that the bottleneck is concentrated in the MLP analog path, while both attention-side linearizations (QKV and projection) collapse structurally. Validation against fabricated organic arrays is deferred to future work. Overall, the framework ranks deployment risks, not a chip-predictive emulator.

\subsection{Outlook}

Next steps include automated profile fitting from laboratory characterization and extended multi-instance HAT schedules. The natural extension of the present behavioral framework is a circuit-aware layer that models spatial IR drop, sneak paths, and temperature-dependent coefficients from array geometry---a scope we have explicitly deferred here to prioritize materials-to-system risk ranking before full hardware closure. Extrapolation to ImageNet-scale deployment is also outside the present evidence base: per-epoch D2D resampling would incur substantially higher training overhead, the 6-bit ADC cliff may shift upward for deeper MLP stacks with finer decision boundaries, and the current results do not address fresh large-scale training from random initialization. A preliminary cross-framework comparison against CrossSim on a deterministic 1\,000-image CIFAR-10 test subset (single-run clean baseline, 3-run Monte Carlo under noise) shows consistent baseline inference (86.2\% ours vs.\ 83.7\% CrossSim at 8-bit ADC) but diverges under noise injection ($81.63 \pm 0.56$\% vs.\ $67.20 \pm 2.67$\% at $\sigma=5\%$, a large qualitative divergence of 14.43~pp at $n=3$, preliminary), highlighting the sensitivity of accuracy predictions to the noise-to-conductance mapping. The subset protocol is disclosed explicitly because of CrossSim throughput constraints; see Supplementary Note~SX.Y.
